\documentclass[11pt,a4paper]{scrartcl}

\usepackage[provide=*,lithuanian]{babel}

\usepackage{../1.\space Vienanariai\space ir\space daugianariai/manostilius_lt}

\usepackage{multicol}

\usepackage{tasks}

\usepackage[left=2.15cm,right=2.15cm,top=2.2cm,bottom=2.2cm,headheight=16pt]{geometry}

\numberwithin{exercise}{section}

\makeatletter

\renewenvironment{exercise}{\@thm{\th@definition}{exercise}{}}{\@endtheorem}

\makeatother

\titlehead{}

\title{Kvadratinio daugianario ekstremumo taškas ir ekstremumai}

\subtitle{2 mokomasis lapas}

\author{Andrius Berniukevičius}

\setlength{\emergencystretch}{2em}

\hypersetup{colorlinks=true,linkcolor=MidnightBlue}

\newcommand{\pagalba}[1]{{\small\itshape Jei kyla sunkumų, žr. \ref{#1} pavyzdį.}}

\begin{document}

\maketitle

\vspace{-1cm}

\section{Ekstremumo taškas, minimumas ir maksimumas}

\subsection{Ekstremumo taškas, minimumas ir maksimumas}

Pirmoje temoje įrodėme: kai $a\ne0$ ir $D=b^2-4ac$,

\[P(x)=ax^2+bx+c=a\left(x+\frac{b}{2a}\right)^2-\frac{D}{4a}.\]

Žymėjimas $P(x)$ nusako daugianarį ir kad kintamasis yra $x$, o pvz. $P(1)$ -- jo reikšmę įrašius $x=1$. 

\begin{definition}\label{ap:ekstremumai}

\vocab{Minimumas} -- mažiausia daugianario reikšmė: visos jo reikšmės yra ne mažesnės už ją ir su bent viena $x$ reikšme ji pasiekiama. \vocab{Maksimumas} -- didžiausia daugianario reikšmė: visos jo reikšmės yra ne didesnės už ją ir su bent viena $x$ reikšme ji pasiekiama. Minimumą ir maksimumą bendrai vadiname \vocab{ekstremumais}.

\end{definition}

\begin{definition}\label{ap:ekstremumo_taskas}

Kvadratinio daugianario $P(x)=a\left(x+\frac{b}{2a}\right)^2-\frac{D}{4a}$ \vocab{ekstremumo taškas} yra tokia kintamojo reikšmė, kuriai reiškinys skliaustuose lygus nuliui. Iš $x+\dfrac{b}{2a}=0$ gauname

\[
\boxed{x=-\frac{b}{2a}}.
\]

Apibendrinant: kvadratinio trinario ekstremumo taškas yra $-\frac{b}{2a}$, o ekstremumo reikšmė -- $-\frac{D}{4a}$.

\end{definition}

\begin{theorem}\label{th:ekstremumai}

Jei $a>0$, tai $P(x)=a\left(x+\frac{b}{2a}\right)^2-\frac{D}{4a}$ mažiausią reikšmę įgyja tada, kai $x=-\frac{b}{2a}$, ir ši reikšmė lygi $-\frac{D}{4a}$.

Jei $a<0$, tas pats reiškinys didžiausią reikšmę įgyja tada, kai
$x=-\frac{b}{2a}$, ir ši reikšmė taip pat lygi $-\frac{D}{4a}$.

\end{theorem}

\begin{proof}

Visada
\[
\left(x+\frac{b}{2a}\right)^2\geq0,
\]
o šis kvadratas lygus nuliui tada ir tik tada, kai
\[
x+\frac{b}{2a}=0,\qquad x=-\frac{b}{2a}.
\]

Jei $a>0$, padauginę nelygybę iš $a$, gauname
\[
a\left(x+\frac{b}{2a}\right)^2\geq0.
\]
Pridėję $-\frac{D}{4a}$, gauname
\[
P(x)=a\left(x+\frac{b}{2a}\right)^2-\frac{D}{4a}
\geq-\frac{D}{4a}.
\]
Lygybė pasiekiama tada ir tik tada, kai
$x=-\frac{b}{2a}$, todėl tai yra mažiausia reikšmė.

Jei $a<0$, padauginę nelygybę iš $a$, apverčiame jos ženklą:
\[
a\left(x+\frac{b}{2a}\right)^2\leq0.
\]
Pridėję $-\frac{D}{4a}$, gauname
\[
P(x)=a\left(x+\frac{b}{2a}\right)^2-\frac{D}{4a}
\leq-\frac{D}{4a}.
\]
Lygybė pasiekiama tada ir tik tada, kai
$x=-\frac{b}{2a}$, todėl tai yra didžiausia reikšmė.

\end{proof}

\Needspace{8\baselineskip}
\begin{example}\label{pv:riba_lygybe}
Užpildykite tarpus taip, kad nelygybės būtų teisingos, ir nurodykite, kada pasiekiama lygybė:
\[
P(x)=3(x+2)^2-7\geq\ldots,
\qquad
Q(x)=-2(x-4)^2+5\leq\ldots.
\]
\end{example}

\begin{solution}
Kadangi $(x+2)^2\geq0$, tai $3(x+2)^2\geq0$. Pridėję $-7$, gauname
\[
P(x)=3(x+2)^2-7\geq-7.
\]
Lygybė pasiekiama tada, kai $(x+2)^2=0$, t. y. kai $x=-2$.

Taip pat $(x-4)^2\geq0$, todėl, padauginę iš neigiamo skaičiaus $-2$, nelygybės ženklą apverčiame:
\[
-2(x-4)^2\leq0.
\]
Pridėję $5$, gauname
\[
Q(x)=-2(x-4)^2+5\leq5.
\]
Lygybė pasiekiama tada, kai $(x-4)^2=0$, t. y. kai $x=4$.

\ats{$P(x)\geq-7$, lygybė pasiekiama, kai $x=-2$; $Q(x)\leq5$, lygybė pasiekiama, kai $x=4$.}
\end{solution}


\Needspace{8\baselineskip}

\begin{example}\label{pv:pavidalas}

Nustatykite ekstremumo tašką ir ekstremumą:

\begin{tasks}[label=\alph*),label-width=1.4em,label-offset=0.1em,item-indent=1.5em](1)

\task $P(x)=2(x+3)^2-4$;

\task $Q(x)=-3(x-2)^2+7$.

\end{tasks}

\end{example}

\begin{solution}\hfill\par

\text{a)} Kvadratas lygus nuliui, kai $x+3=0$, taigi ekstremumo taškas yra $x=-3$. Kadangi $a=2>0$, tai $P(x)\geq-4$ ir lygybė pasiekiama, kai $x=-3$. \emph{Minimumas lygus $-4$; maksimumo nėra.}



\text{b)} Kvadratas lygus nuliui, kai $x-2=0$, taigi ekstremumo taškas yra $x=2$. Kadangi $a=-3<0$, tai $Q(x)\leq7$ ir lygybė pasiekiama, kai $x=2$. \emph{Maksimumas lygus $7$; minimumo nėra.}

\ats{\begin{tasks}[label=\alph*),label-width=1.4em,label-offset=0.1em,item-indent=1.5em](1)

\task Ekstremumo taškas $x=-3$, minimumas $-4$; maksimumo nėra.

\task Ekstremumo taškas $x=2$, maksimumas $7$; minimumo nėra.

\end{tasks}}

\end{solution}

\Needspace{8\baselineskip}

\begin{example}\label{pv:trinaris}

Raskite ekstremumo tašką ir ekstremumą:

\begin{tasks}[label=\alph*),label-width=1.4em,label-offset=0.1em,item-indent=1.5em](1)

\task $P(x)=2x^2-6x+7$;

\task $Q(x)=-x^2-4x+1$.

\end{tasks}

\end{example}

\begin{solution}\hfill\break

\textbf{a)} $a=2$, $b=-6$, $c=7$, todėl

\[D=36-56=-20,\quad x=-\frac{-6}{4}=\frac32,
\quad P\left(\frac32\right)=-\frac{-20}{8}=\frac52.\]

Kadangi $a>0$, tai minimumas. Galime patikrinti išskyrę kvadratą: $$P(x)=2\left(x-\frac{3}{2}\right)^2+\frac{5}{2}.$$ 

\emph{Ekstremumo taškas $x=\frac32$, minimumas $\frac52$; maksimumo nėra.}



\textbf{b)} $a=-1$, $b=-4$, $c=1$, todėl

\[D=16+4=20,\quad x=-\frac{-4}{-2}=-2,
\quad Q(-2)=-\frac{20}{-4}=5.\]

Kadangi $a<0$, tai maksimumas. Iš tiesų $Q(x)=-(x+2)^2+5$. 

\emph{Ekstremumo taškas $x=-2$, maksimumas $5$; minimumo nėra.}

\ats{\begin{tasks}[label=\alph*),label-width=1.4em,label-offset=0.1em,item-indent=1.5em](1)

\task Ekstremumo taškas $x=\frac32$, minimumas $\frac52$; maksimumo nėra.

\task Ekstremumo taškas $x=-2$, maksimumas $5$; minimumo nėra.

\end{tasks}}

\end{solution}

\Needspace{8\baselineskip}

\begin{example}\label{pv:atkurti}

Raskite kvadratinį daugianarį $P(x)$, jei jo ekstremumo taškas yra $x=-1$, ten pasiekiamas minimumas $-3$, o $P(1)=5$.

\end{example}

\begin{solution}

Ekstremumo taškas ir minimumo reikšmė nurodo pavidalą $P(x)=a(x+1)^2-3$, tačiau dar nežinome $a$. Sąlyga $P(1)=5$ reiškia, kad, įrašę $x=1$, turime gauti $5$:

\[a(1+1)^2-3=5,\qquad 4a=8,\qquad a=2.\]

Koeficientas teigiamas, todėl nurodytas ekstremumas tikrai yra minimumas. Gauname

\[\boxed{P(x)=2(x+1)^2-3=2x^2+4x-1.}\]

Patikra: $P(1)=8-3=5$; kvadratas lygus nuliui, kai $x=-1$, o ten daugianario reikšmė lygi $-3$.






\ats{$P(x)=2(x+1)^2-3=2x^2+4x-1$.}

\end{solution}

\subsection{Daugianario reikšmių palyginimas}

Jei daugianaris užrašytas pavidalu $P(x)=a(x-h)^2+k$, įrašius skaičių $u$ vietoje $x$ gaunama $P(u)=a(u-h)^2+k$. Lyginant dvi reikšmes, pakanka palyginti atitinkamus kvadratus: kai $a>0$, didesnis kvadratas duoda didesnę reikšmę, o kai $a<0$ -- mažesnę. Jei kvadratai lygūs, daugianario reikšmės taip pat lygios.

\Needspace{8\baselineskip}
\begin{example}\label{pv:palyginti_reiksmes}
Duota $P(x)=2(x-3)^2-4$ ir $Q(x)=-2(x+1)^2+3$. Įrašykite ženklą $<$, $>$ arba $=$, neskaičiuodami atskirų daugianario reikšmių:
\[
\begin{aligned}
P(4)&\ \ldots\ P(6),\\
P(1)&\ \ldots\ P(5),\\
P(-1)&\ \ldots\ P(4),\\
Q(-2)&\ \ldots\ Q(2),\\
Q(-3)&\ \ldots\ Q(1),\\
Q(4)&\ \ldots\ Q(0).
\end{aligned}
\]
\end{example}

\begin{solution}
Kai $P(x)=2(x-3)^2-4$, kvadratas padaugintas iš teigiamo skaičiaus $2$. Todėl didesnis kvadratas duoda didesnę reikšmę:
\begin{itemize}
\item Kai $x=4$ ir $x=6$, gauname atitinkamai $(4-3)^2=1$ ir $(6-3)^2=9$. Kadangi $1<9$, tai $P(4)<P(6)$.
\item Kai $x=1$ ir $x=5$, gauname atitinkamai $(1-3)^2=4$ ir $(5-3)^2=4$. Kvadratai lygūs, todėl $P(1)=P(5)$.
\item Kai $x=-1$ ir $x=4$, gauname atitinkamai $(-1-3)^2=16$ ir $(4-3)^2=1$. Kadangi $16>1$, tai $P(-1)>P(4)$.
\end{itemize}

Kai $Q(x)=-2(x+1)^2+3$, kvadratas padaugintas iš neigiamo skaičiaus $-2$. Didesnį kvadratą padauginus iš neigiamo skaičiaus gaunama mažesnė reikšmė:
\begin{itemize}
\item Kai $x=-2$ ir $x=2$, gauname atitinkamai $(-2+1)^2=1$ ir $(2+1)^2=9$. Kadangi $1<9$, tai $Q(-2)>Q(2)$.
\item Kai $x=-3$ ir $x=1$, gauname atitinkamai $(-3+1)^2=4$ ir $(1+1)^2=4$. Kvadratai lygūs, todėl $Q(-3)=Q(1)$.
\item Kai $x=4$ ir $x=0$, gauname atitinkamai $(4+1)^2=25$ ir $(0+1)^2=1$. Kadangi $25>1$, tai $Q(4)<Q(0)$.
\end{itemize}
\ats{$P(4)<P(6)$, $P(1)=P(5)$, $P(-1)>P(4)$; $Q(-2)>Q(2)$, $Q(-3)=Q(1)$, $Q(4)<Q(0)$.}
\end{solution}

\section{Sunkesnių uždavinių sprendimas}

\Needspace{8\baselineskip}
\begin{example}\label{pv:sunk_diskriminantas_vienodos}
Raskite kvadratinį trinarį $P(x)$, jei jo diskriminantas lygus
$-128$, ekstremumas lygus $4$, o $P(-1)=P(5)$.
\end{example}

\begin{solution}
Ekstremumo reikšmė lygi
\[
-\frac{D}{4a}.
\]
Įrašę duotus skaičius, gauname
\[
4=-\frac{-128}{4a}=\frac{32}{a},
\qquad a=8.
\]
Gavome $a=8>0$, todėl ekstremumas yra minimumas.

Pažymėkime $h=\dfrac{b}{2a}$. Tuomet trinarį galime užrašyti pavidalu
\[
P(x)=8(x+h)^2+4,
\]
o jo ekstremumo taškas yra $x=-h$.
Sąlyga $P(-1)=P(5)$ reiškia, kad
\[
8(-1+h)^2+4=8(5+h)^2+4.
\]
Atėmę iš abiejų pusių $4$ ir padaliję iš $8$, gauname
\[
(-1+h)^2=(5+h)^2.
\]
Išskleidžiame kvadratus:
\[
1-2h+h^2=25+10h+h^2.
\]
Todėl
\[
12h=-24,\qquad h=-2.
\]
Vadinasi, ekstremumo taškas yra $x=-h=2$ ir
\[
P(x)=8(x-2)^2+4=8x^2-32x+36.
\]

Patikriname duomenis:
\[
D=(-32)^2-4\cdot8\cdot36=-128,
\]
\[
P(-1)=76,\qquad P(5)=76.
\]
Be to, $8(x-2)^2+4\ge4$, o lygybė pasiekiama, kai $x=2$,
taigi minimumas tikrai lygus $4$.

\ats{$P(x)=8x^2-32x+36$; ekstremumas yra minimumas.}
\end{solution}

\Needspace{8\baselineskip}
\begin{example}\label{pv:sunk_du_variantai}
Raskite visus kvadratinius trinarius $P(x)$, kurių diskriminantas
lygus $-72$, ekstremumas lygus $-3$, o $P(1)=-27$.
\end{example}

\begin{solution}
Ekstremumo reikšmė lygi $-\dfrac{D}{4a}$, todėl
\[
-3=-\frac{-72}{4a}=\frac{18}{a},
\qquad a=-6.
\]
Gavome $a=-6<0$, todėl ekstremumas yra maksimumas.

Pažymėkime $h=\dfrac{b}{2a}$. Tuomet trinarį galime užrašyti pavidalu
\[
P(x)=-6(x+h)^2-3,
\]
o jo ekstremumo taškas yra $x=-h$.
Pasinaudojame sąlyga $P(1)=-27$:
\[
-6(1+h)^2-3=-27,
\]
\[
6(1+h)^2=24,
\qquad (1+h)^2=4.
\]
Skaičiaus kvadratas lygus $4$ tada ir tik tada, kai pats skaičius
lygus $2$ arba $-2$. Todėl būtina išnagrinėti abu atvejus:
\[
1+h=2 \quad\text{arba}\quad 1+h=-2.
\]
Gauname
\[
h=1 \quad\text{arba}\quad h=-3.
\]

Jei $h=1$, tai
\[
P(x)=-6(x+1)^2-3=-6x^2-12x-9.
\]
Jei $h=-3$, tai
\[
P(x)=-6(x-3)^2-3=-6x^2+36x-57.
\]



\ats{$P(x)=-6x^2-12x-9$ arba $P(x)=-6x^2+36x-57$; ekstremumas yra maksimumas.}
\end{solution}

\Needspace{8\baselineskip}
\begin{example}\label{pv:sunk_koeficientu_rysiai}
Kvadratinio trinario $P(x)=ax^2+bx+c$ koeficientai tenkina sąlygas
\[
a+c=13,\qquad ac=36,\qquad P(1)=7.
\]
Raskite visus tokius trinarius ir jų ekstremumo taškus.
\end{example}

\begin{solution}
Įrašę $x=1$ į trinario išraišką, gauname
\[
P(1)=a+b+c.
\]
Todėl
\[
a+b+c=7.
\]
Kadangi $a+c=13$, tai
\[
13+b=7,\qquad b=-6.
\]

Lieka rasti $a$ ir $c$. Iš $a+c=13$ išreiškiame
\[
c=13-a.
\]
Įrašę į sąlygą $ac=36$, gauname
\[
a(13-a)=36,
\]
\[
a^2-13a+36=0,
\]
\[
(a-4)(a-9)=0.
\]
Vadinasi, $a=4$ arba $a=9$.

\fbox{\textbf{1 atvejis:}} $a=4$. Tuomet $c=9$, todėl
\[
P(x)=4x^2-6x+9.
\]
Kadangi $a>0$, tai ekstremumo taškas yra
\[
x=-\frac{-6}{2\cdot4}=\frac34.
\]

\fbox{\textbf{2 atvejis:}} $a=9$. Tuomet $c=4$, todėl
\[
P(x)=9x^2-6x+4.
\]
Kadangi ir šiuo atveju $a>0$, ekstremumo taškas yra
\[
x=-\frac{-6}{2\cdot9}=\frac13.
\]

Patikriname: abiem atvejais $a+c=13$, $ac=36$, o
\[
P(1)=a+b+c=13-6=7.
\]

\ats{$P(x)=4x^2-6x+9$, ekstremumo taškas $x=\frac34$;
arba $P(x)=9x^2-6x+4$, ekstremumo taškas $x=\frac13$.}
\end{solution}

\Needspace{8\baselineskip}
\begin{example}\label{pv:sunk_trys_reiksmes}
Žinoma, kad kvadratinio trinario $P(x)$ reikšmės yra
\[
P(0)=8,\qquad P(1)=-1,\qquad P(2)=-4.
\]
Raskite trinario ekstremumo tašką ir nustatykite,
ar jame pasiekiamas minimumas, ar maksimumas.
\end{example}

\begin{solution}
Užrašome
\[
P(x)=ax^2+bx+c.
\]
Kadangi $P(0)=c$, iš pirmosios sąlygos iš karto gauname
\[
c=8.
\]
Įrašę $x=1$ ir $x=2$, gauname dar dvi lygybes:
\[
a+b+8=-1,
\qquad
4a+2b+8=-4.
\]
Sutvarkome:
\[
\begin{cases}
a+b=-9,\\
4a+2b=-12.
\end{cases}
\]
Pirmąją lygybę padauginame iš $2$:
\[
2a+2b=-18.
\]
Atimame ją iš antrosios lygybės:
\[
(4a+2b)-(2a+2b)=-12-(-18),
\]
\[
2a=6,\qquad a=3.
\]
Tuomet
\[
3+b=-9,\qquad b=-12.
\]
Taigi
\[
P(x)=3x^2-12x+8.
\]

Kadangi $a=3>0$, trinaris turi minimumą. Ekstremumo taškas:
\[
x=-\frac{b}{2a}
=-\frac{-12}{2\cdot3}=2.
\]
Galime patikrinti išskyrę pilnąjį kvadratą:
\[
P(x)=3(x^2-4x)+8
=3\bigl((x-2)^2-4\bigr)+8
=3(x-2)^2-4.
\]
Kvadratas $(x-2)^2$ yra neneigiamas ir lygus nuliui tik kai $x=2$.
Todėl minimumas pasiekiamas būtent šiame taške.

\ats{Ekstremumo taškas $x=2$; minimumas.}
\end{solution}

\section{Formuojame automatizmą}

\begin{exercise}\label{uz:riba_lygybe}
Užpildykite tarpus taip, kad nelygybės būtų teisingos, ir nurodykite, kada pasiekiama lygybė.

\par\pagalba{pv:riba_lygybe}

\begin{tasks}[label=(\arabic*),label-width=2em](1)

\task $P(x)=3(x+2)^2-7\geq\ldots$; lygybė pasiekiama, kai $x=\ldots$.

\task $Q(x)=-2(x-4)^2+5\leq\ldots$; lygybė pasiekiama, kai $x=\ldots$.

\task $R(x)=4(x-1)^2+3\geq\ldots$; lygybė pasiekiama, kai $x=\ldots$.

\task $S(x)=-5(x+3)^2-2\leq\ldots$; lygybė pasiekiama, kai $x=\ldots$.

\task $T(x)=\frac12(x-6)^2-4\geq\ldots$; lygybė pasiekiama, kai $x=\ldots$.

\task $U(x)=-\frac13(x+1)^2+8\leq\ldots$; lygybė pasiekiama, kai $x=\ldots$.

\task $V(x)=7(x+2)^2\geq\ldots$; lygybė pasiekiama, kai $x=\ldots$.

\task $W(x)=-4(x-5)^2+1\leq\ldots$; lygybė pasiekiama, kai $x=\ldots$.

\task $F(x)=2(x-3)^2-9\geq\ldots$; lygybė pasiekiama, kai $x=\ldots$.

\task $G(x)=-6(x+1)^2-4\leq\ldots$; lygybė pasiekiama, kai $x=\ldots$.

\end{tasks}
\end{exercise}
\newpage
\begin{exercise}\label{uz:pavidalas}

Nurodykite daugianario ekstremumo tašką ir ekstremumo reikšmę. Parašykite, ar tai minimumas, ar maksimumas, ir nurodykite, kurios rūšies ekstremumo nėra. Pavyzdžiui: „Ekstremumo taškas $x=\ldots$, minimumas lygus $\ldots$, maksimumo nėra.“

\par\pagalba{pv:pavidalas}

\begin{tasks}[label=(\arabic*),label-width=2em](2)

\task $P(x)=2\left(x-3\right)^2-5$.

\task $P(x)=-\left(x+2\right)^2+8$.

\task $P(x)=3\left(x+4\right)^2+1$.

\task $P(x)=-2\left(x-5\right)^2-3$.

\task $P(x)=\frac{1}{2}\left(x+1\right)^2-6$.

\task $P(x)=-4\left(x+3\right)^2$.

\task $P(x)=\left(x-\frac{3}{2}\right)^2+\frac{7}{4}$.

\task $P(x)=-3\left(x+\frac{1}{2}\right)^2+\frac{5}{2}$.

\task $P(x)=5x^2-7$.

\task $P(x)=-x^2-4$.

\end{tasks}

\end{exercise}

\begin{exercise}\label{uz:palyginti_reiksmes}
Įrašykite ženklą $<$, $>$ arba $=$, neskaičiuodami atskirų daugianario reikšmių.

\par\pagalba{pv:palyginti_reiksmes}

\begin{tasks}[label=(\arabic*),label-width=2em](1)

\task $P(x)=2(x-3)^2-4$: \quad $P(4)\ \ldots\ P(6)$, \quad $P(1)\ \ldots\ P(5)$, \quad $P(-1)\ \ldots\ P(4)$.

\task $Q(x)=-3(x+2)^2+5$: \quad $Q(-1)\ \ldots\ Q(1)$, \quad $Q(-5)\ \ldots\ Q(1)$, \quad $Q(2)\ \ldots\ Q(-2)$.

\task $R(x)=\frac12(x-1)^2+2$: \quad $R(0)\ \ldots\ R(3)$, \quad $R(-1)\ \ldots\ R(3)$, \quad $R(4)\ \ldots\ R(0)$.

\task $S(x)=-2(x-4)^2+1$: \quad $S(3)\ \ldots\ S(6)$, \quad $S(2)\ \ldots\ S(6)$, \quad $S(0)\ \ldots\ S(5)$.

\task $T(x)=5(x+1)^2-3$: \quad $T(-2)\ \ldots\ T(1)$, \quad $T(-4)\ \ldots\ T(2)$, \quad $T(3)\ \ldots\ T(0)$.

\task $U(x)=-\frac14(x-2)^2+7$: \quad $U(1)\ \ldots\ U(4)$, \quad $U(-1)\ \ldots\ U(5)$, \quad $U(6)\ \ldots\ U(2)$.

\task $V(x)=3(x-5)^2$: \quad $V(4)\ \ldots\ V(7)$, \quad $V(3)\ \ldots\ V(7)$, \quad $V(0)\ \ldots\ V(6)$.

\task $W(x)=-6(x+3)^2-2$: \quad $W(-2)\ \ldots\ W(0)$, \quad $W(-5)\ \ldots\ W(-1)$, \quad $W(2)\ \ldots\ W(-3)$.

\task $F(x)=\frac32(x-2)^2-8$: \quad $F(1)\ \ldots\ F(4)$, \quad $F(-1)\ \ldots\ F(5)$, \quad $F(6)\ \ldots\ F(0)$.

\task $G(x)=-2(x-1)^2+6$: \quad $G(0)\ \ldots\ G(3)$, \quad $G(-2)\ \ldots\ G(4)$, \quad $G(5)\ \ldots\ G(1)$.

\end{tasks}
\end{exercise}

\begin{exercise}\label{uz:trinaris}

Raskite daugianario ekstremumo tašką ir ekstremumo reikšmę. Nurodykite, ar tai minimumas, ar maksimumas, ir kuri ekstremumo rūšis neegzistuoja.

\par\pagalba{pv:trinaris}

\begin{tasks}[label=(\arabic*),label-width=2em](2)

\task $P(x)=x^2-4x+19$.

\task $P(x)=x^2+6x+5$.

\task $P(x)=2x^2-8x+3$.

\task $P(x)=-x^2+10x-7$.

\task $P(x)=3x^2+6x+8$.

\task $P(x)=-2x^2-4x+9$.

\task $P(x)=4x^2+4x+3$.

\task $P(x)=-3x^2+3x+2$.

\task $P(x)=2x^2-5$.

\task $P(x)=-2x^2+6x$.

\end{tasks}

\end{exercise}

\newpage

\begin{exercise}\label{uz:atkurti}

Kiekvienu atveju raskite kvadratinį daugianarį $P(x)=ax^2+bx+c$, kai duotas jo ekstremumo taškas, ten pasiekiamas minimumas arba maksimumas ir papildoma daugianario reikšmė.

\par\pagalba{pv:atkurti}

\begin{tasks}[label=(\arabic*), label-width=1.8em](2)

\task Ekstremumo taškas $x=2$, minimumas $3$, $P(1)=5$.

\task Ekstremumo taškas $x=-1$, maksimumas $9$, $P(1)=5$.

\task Ekstremumo taškas $x=3$, minimumas $1$, $P(0)=10$.

\task Ekstremumo taškas $x=0$, maksimumas $8$, $P(2)=-4$.

\task Ekstremumo taškas $x=-2$, minimumas $\frac{1}{2}$, $P(0)=\frac{5}{2}$.

\task Ekstremumo taškas $x=4$, maksimumas $23$, $P(2)=15$.

\task Ekstremumo taškas $x=-3$, minimumas $-11$, $P(-1)=-7$.

\task Ekstremumo taškas $x=-1$, maksimumas $7$, $P(0)=\frac{13}{2}$.

\task Ekstremumo taškas $x=3$, minimumas $-3$, $P(2)=-1$.

\task Ekstremumo taškas $x=2$, maksimumas $9$, $P(0)=-7$.

\end{tasks}

\end{exercise}

\section{Sudėtingesni uždaviniai}

\begin{exercise} Kiekvienu atveju raskite trinarį $P(x) = ax ^{2} +bx+c$.

\begin{tasks}[label=(\arabic*),label-width=1.8em,label-offset=0.1em,item-indent=2em](1)

\task Trinario diskriminantas lygus $-64$, minimumas lygus $4$ ir pasiekiamas, kai $x=3$.

\task Trinario diskriminantas lygus $-144$, minimumas lygus $6$, o $P(0)=P(4)$.

\task Trinario diskriminantas lygus $-100$, ekstremumo taškas yra $x=2$, o minimumas lygus $5$.

\task Trinario diskriminantas lygus $-81$, minimumas lygus $9$, o $P(1)=P(3)$.

\task Trinario diskriminantas lygus $64$, ekstremumas lygus $-4$ ir pasiekiamas, kai $x=3$.

\task Trinario diskriminantas lygus $144$, ekstremumas lygus $-6$, o $P(0)=P(4)$. Raskite trinarį ir nustatykite, ar ekstremumas yra minimumas, ar maksimumas.

\task Trinario diskriminantas lygus $100$, ekstremumas lygus $5$, o $P(0)=5$. Raskite trinarį ir nustatykite, ar ekstremumas yra minimumas, ar maksimumas.

\task Trinario ekstremumas lygus $-8$, diskriminantas lygus $144$, o $P(0)=10$.

\task Trinario diskriminantas lygus $144$, ekstremumas lygus $-9$, o $P(3)=0$.

\task Trinario diskriminantas lygus $64$, ekstremumas lygus $4$, o $P(2)=-5$.

\task Trinario diskriminantas lygus $81$, ekstremumas lygus $9$, o $P(1)=0$.

\task Trinario diskriminantas lygus $100$, ekstremumas lygus $-5$, o $P(1)=P(5)+20$.

\task Trinario koeficientai tenkina lygybes $a+c=10$, $ac=21$, $P(1)=6$.

\task Trinario diskriminantas lygus $36$, ekstremumas lygus $-3$, o $P(1)=9$.

\task Žinoma, kad $P(1)-P(-1)=8$, $ac=3$, $P(2)=15$.

\task Trinario koeficientai tenkina lygybes $\sqrt{a}+\sqrt{c}=5$, $ac=36$, $P(1)=0$.

\task Žinoma, kad
\[
\begin{cases}
b-\sqrt{ac}=2,\\
b+\sqrt{ac}=8,
\end{cases}
\]
o trinario mažiausioji reikšmė lygi $\frac{11}{4}$.

\task Trinario diskriminantas lygus $100$, ekstremumas lygus $-5$, o $P(1)+P(3)=10$.

\task Trinario ekstremumas lygus $-9$ ir pasiekiamas, kai $x=2$. Be to, $P(0)\cdot P(4)=25$.

\task Trinario koeficientas $a>0$, diskriminantas lygus $80$, o $P(1)=P(5)$, $P(0)=8$.

\task Trinario diskriminantas lygus nuliui. Be to, $P(1)=4$, $P(3)=16$.

\task Abi trinario šaknys yra realiosios, nelygios nuliui, o jų sandauga lygi $1$. Be to, $P(1)=-1$, $D=9$.

\task Žinoma, kad
\[
\begin{cases}
P(1)+P(-1)=20,\\
P(1)\cdot P(-1)=64,\\
ac=21.
\end{cases}
\]

\task Trinario koeficientai tenkina lygybes $b^2=ac$, $P(1)=6$, $P(-1)=2$.

\end{tasks}
\end{exercise}

\begin{exercise}\label{uz:taskas_vienodos_reiksmes}
Žinoma, kad trinario $P(x)=ax^2+bx+c$ koeficientas $a>0$, o $P(-2)=P(8)$. Raskite trinario ekstremumo tašką.
\end{exercise}

\begin{exercise}\label{uz:taskas_saknys}
Trinario $P(x)=ax^2+bx+c$ koeficientas $a<0$. Jo realiosios šaknys $x_1$, $x_2$ tenkina sąlygas $x_1+x_2=10$, $x_1x_2=16$. Raskite trinario ekstremumo tašką.
\end{exercise}

\begin{exercise}\label{uz:taskas_trys_reiksmes}
Žinoma, kad kvadratinio trinario $P(x)$ reikšmės yra $P(0)=7$, $P(1)=2$, $P(2)=1$. Raskite trinario ekstremumo tašką ir nustatykite, ar jame pasiekiamas minimumas, ar maksimumas.
\end{exercise}

\begin{exercise}\label{uz:taskas_diskriminantas}
Trinario diskriminantas lygus $36$, minimumas lygus $-3$, o $P(1)=9$. Raskite visas galimas trinario ekstremumo taško reikšmes.
\end{exercise}

\begin{exercise}\label{uz:taskas_koeficientai}
Trinario $P(x)=ax^2+bx+c$ koeficientai tenkina sąlygas $a+c=10$, $ac=21$, $P(1)=6$. Raskite visas galimas trinario ekstremumo taško reikšmes ir kiekvienu atveju nustatykite, ar ten pasiekiamas minimumas, ar maksimumas.
\end{exercise}

\newpage


\section{Atsakymai}


\medskip
\noindent\textbf{\ref{uz:riba_lygybe} užduotis.}

\begin{tasks}[label=(\arabic*),label-width=2em](1)
\task $P(x)\geq-7$; lygybė pasiekiama, kai $x=-2$.
\task $Q(x)\leq5$; lygybė pasiekiama, kai $x=4$.
\task $R(x)\geq3$; lygybė pasiekiama, kai $x=1$.
\task $S(x)\leq-2$; lygybė pasiekiama, kai $x=-3$.
\task $T(x)\geq-4$; lygybė pasiekiama, kai $x=6$.
\task $U(x)\leq8$; lygybė pasiekiama, kai $x=-1$.
\task $V(x)\geq0$; lygybė pasiekiama, kai $x=-2$.
\task $W(x)\leq1$; lygybė pasiekiama, kai $x=5$.
\task $F(x)\geq-9$; lygybė pasiekiama, kai $x=3$.
\task $G(x)\leq-4$; lygybė pasiekiama, kai $x=-1$.
\end{tasks}

\medskip
\noindent\textbf{\ref{uz:pavidalas} užduotis.}

\begin{tasks}[label=(\arabic*),label-width=2em](2)
\task $x=3$; minimumas $-5$.
\task $x=-2$; maksimumas $8$.
\task $x=-4$; minimumas $1$.
\task $x=5$; maksimumas $-3$.
\task $x=-1$; minimumas $-6$.
\task $x=-3$; maksimumas $0$.
\task $x=\frac32$; minimumas $\frac74$.
\task $x=-\frac12$; maksimumas $\frac52$.
\task $x=0$; minimumas $-7$.
\task $x=0$; maksimumas $-4$.
\end{tasks}

\medskip
\noindent\textbf{\ref{uz:palyginti_reiksmes} užduotis.}

\begin{tasks}[label=(\arabic*),label-width=2em](1)
\task $P(4)<P(6)$, $P(1)=P(5)$, $P(-1)>P(4)$.
\task $Q(-1)>Q(1)$, $Q(-5)=Q(1)$, $Q(2)<Q(-2)$.
\task $R(0)<R(3)$, $R(-1)=R(3)$, $R(4)>R(0)$.
\task $S(3)>S(6)$, $S(2)=S(6)$, $S(0)<S(5)$.
\task $T(-2)<T(1)$, $T(-4)=T(2)$, $T(3)>T(0)$.
\task $U(1)>U(4)$, $U(-1)=U(5)$, $U(6)<U(2)$.
\task $V(4)<V(7)$, $V(3)=V(7)$, $V(0)>V(6)$.
\task $W(-2)>W(0)$, $W(-5)=W(-1)$, $W(2)<W(-3)$.
\task $F(1)<F(4)$, $F(-1)=F(5)$, $F(6)>F(0)$.
\task $G(0)>G(3)$, $G(-2)=G(4)$, $G(5)<G(1)$.
\end{tasks}

\medskip
\noindent\textbf{\ref{uz:trinaris} užduotis.}

\begin{tasks}[label=(\arabic*),label-width=2em](2)
\task $x=2$; minimumas $15$.
\task $x=-3$; minimumas $-4$.
\task $x=2$; minimumas $-5$.
\task $x=5$; maksimumas $18$.
\task $x=-1$; minimumas $5$.
\task $x=-1$; maksimumas $11$.
\task $x=-\frac12$; minimumas $2$.
\task $x=\frac12$; maksimumas $\frac{11}{4}$.
\task $x=0$; minimumas $-5$.
\task $x=\frac32$; maksimumas $\frac92$.
\end{tasks}

\medskip
\noindent\textbf{\ref{uz:atkurti} užduotis.}

\begin{tasks}[label=(\arabic*),label-width=2em](2)
\task $P(x)=2x^2-8x+11$.
\task $P(x)=-x^2-2x+8$.
\task $P(x)=x^2-6x+10$.
\task $P(x)=-3x^2+8$.
\task $P(x)=\frac12x^2+2x+\frac52$.
\task $P(x)=-2x^2+16x-9$.
\task $P(x)=x^2+6x-2$.
\task $P(x)=-\frac12x^2-x+\frac{13}{2}$.
\task $P(x)=2x^2-12x+15$.
\task $P(x)=-4x^2+16x-7$.
\end{tasks}

\medskip
\noindent\textbf{4.1 užduotis.}

\begin{tasks}[label=(\arabic*),label-width=2em](1)
\task $P(x)=4x^2-24x+40$.

\task $P(x)=6x^2-24x+30$.

\task $P(x)=5x^2-20x+25$.

\task $P(x)=\frac94x^2-9x+18$.

\task $P(x)=4x^2-24x+32$.

\task $P(x)=6x^2-24x+18$; ekstremumas yra minimumas.

\task $P(x)=-5x^2+5$; ekstremumas yra maksimumas.

\task $P(x)=\frac92x^2-18x+10$
arba $P(x)=\frac92x^2+18x+10$.

\task $P(x)=4x^2-12x$
arba $P(x)=4x^2-36x+72$.

\task $P(x)=-4x^2+4x+3$
arba $P(x)=-4x^2+28x-45$.

\task $P(x)=-\frac94x^2-\frac92x+\frac{27}{4}$
arba $P(x)=-\frac94x^2+\frac{27}{2}x-\frac{45}{4}$.

\task $P(x)=5x^2-35x+\frac{225}{4}$.

\task $P(x)=3x^2-4x+7$
arba $P(x)=7x^2-4x+3$.

\task $P(x)=3x^2+6x$
arba $P(x)=3x^2-18x+24$.

\task $P(x)=x^2+4x+3$
arba $P(x)=\frac34x^2+4x+4$.

\task $P(x)=4x^2-13x+9$
arba $P(x)=9x^2-13x+4$.

\task $P(x)=x^2+5x+9$.

\task $P(x)=5x^2-10x$
arba $P(x)=5x^2-30x+40$.

\task $P(x)=\frac72x^2-14x+5$
arba $P(x)=x^2-4x-5$.

\task $P(x)=2x^2-12x+8$.

\task $P(x)=x^2+2x+1$
arba $P(x)=9x^2-30x+25$.

\task $P(x)=2x^2-5x+2$.

\task $P(x)=3x^2-6x+7$,
$P(x)=3x^2+6x+7$,
$P(x)=7x^2-6x+3$
arba $P(x)=7x^2+6x+3$.

\task $P(x)=2x^2+2x+2$.
\end{tasks}

\medskip
\noindent\textbf{\ref{uz:taskas_vienodos_reiksmes} užduotis.}
Ekstremumo taškas $x=3$.

\medskip
\noindent\textbf{\ref{uz:taskas_saknys} užduotis.}
Ekstremumo taškas $x=5$.

\medskip
\noindent\textbf{\ref{uz:taskas_trys_reiksmes} užduotis.}
Ekstremumo taškas $x=\frac74$; jame pasiekiamas minimumas.

\medskip
\noindent\textbf{\ref{uz:taskas_diskriminantas} užduotis.}
Ekstremumo taškas $x=-1$ arba $x=3$.

\medskip
\noindent\textbf{\ref{uz:taskas_koeficientai} užduotis.}
Ekstremumo taškas $x=\frac23$ arba $x=\frac27$;
abiem atvejais jame pasiekiamas minimumas.
\end{document}